\documentclass[12pt]{article}
\usepackage{amsmath, amssymb}
\usepackage{geometry}
\usepackage{hyperref}
\geometry{margin=1in}

\title{On the Nature of Mass:\\
Complexity Generation and Spacetime Curvature}
\author{Paul A. I. Reilly}
\date{[DRAFT v3.0 --- not for circulation]}

\begin{document}

\maketitle

%% ============================================================
\begin{abstract}
%% ============================================================

We show that the exterior Schwarzschild metric can be derived from one
central hypothesis and four supporting assumptions, without invoking
Einstein's equations.
The central hypothesis is that mass is a rate of irreversible information
generation: an isolated mass $M$ of pressureless dust produces classical
information at rate $dC/d\tau = \alpha Mc^2/\hbar$, where $\alpha$ is a
single free parameter.
Extensions to matter with pressure, shear, or momentum flux are addressed
in the companion paper (Paper~II).
Four assumptions govern how this information is geometrically accommodated:
(1) information propagates at the speed of light; (2) the vacuum near
matter operates at holographic saturation; (3) each nat of information
requires a minimum of four Planck areas (Bekenstein); and (4) nature is
parsimonious.
We first fix $\alpha = 2\pi$ uniquely, without reference to the
Schwarzschild metric, by requiring that the information in transit between
the source and any bounding sphere saturate the Bekenstein bound
exactly (Section~2).
With $\alpha$ so fixed, the radial metric component
$g_{rr} = (1 - r_S/r)^{-1}$ follows algebraically as a prediction;
matching to the known Schwarzschild solution provides confirmation rather
than definition (Section~3).
We further present a well-motivated argument for the temporal component
$g_{tt}$, identifying one step as a conjecture open to community
scrutiny; an independent derivation of $g_{tt}$ appears in the companion
paper via the Einstein equations, themselves derived from Screen Theory
premises.
We then show that the same framework, applied to a spherically symmetric
charged mass, recovers the Reissner--Nordstr\"om metric without invoking
Maxwell's equations or the electromagnetic stress-energy tensor
(Section~4).
The structural unity of the Schwarzschild and Reissner--Nordstr\"om
derivations within a single two-parameter framework provides evidence
against an accidental-concordance reading of either result.
A cosmological motivation for the holographic saturation assumption
is developed in Appendix~\ref{app:A2}.

\end{abstract}

%% ============================================================
\section{Introduction}
%% ============================================================

\subsection{The Central Claim}

Mass is ordinarily defined operationally: it measures resistance to
acceleration and determines the strength of gravitational interaction.
In this paper we propose a more fundamental characterization.

\emph{Mass is the rate at which a system generates irreversible classical
information.}

Physics knows how mass behaves: it resists acceleration, sources
gravity, couples to the Higgs field.
It does not know what mass \emph{is}.
The Higgs mechanism specifies the origin of the mass parameter in the
Standard Model Lagrangian without addressing its geometric meaning.
General relativity treats mass-energy as a source term in the field
equations without asking what it is doing to produce curvature.
The present paper proposes that curvature is not a response to mass but
the geometric record of mass's ongoing activity: the irreversible
narrowing of the universe's set of possible futures, at a rate
proportional to the mass-energy present.

Spacetime curvature is the geometric accommodation required by this
ongoing process. More precisely, we propose the hypothesis
\begin{equation}
  \frac{dC}{d\tau} = \alpha \frac{Mc^2}{\hbar},
  \label{eq:hypothesis}
\end{equation}
where $C$ denotes the amount of irreducible classical information
generated by an isolated mass $M$, and $\alpha$ is a single free
parameter fixed in Section~2 by the Bekenstein bound, without reference
to the Schwarzschild metric.
In the derivation below we show that the radial component of the exterior
Schwarzschild metric follows directly from this hypothesis together with
four physical assumptions governing how information is geometrically
accommodated; the Schwarzschild result then constitutes a prediction
confirmed by matching.
Hypothesis~\eqref{eq:hypothesis} is stated for pressureless dust;
extensions to matter with pressure, shear, or momentum flux are more
complex and are addressed in the companion paper (Paper~II).
Pressureless dust isolates the mass-energy density contribution to the
stress-energy tensor cleanly; pressure and momentum flux introduce
additional terms in $T_{\mu\nu}$ that modify the information generation
rate in ways requiring the full treatment of Paper~II.

In this framework classical information is not conserved but increases
monotonically through irreversible physical processes.
The accumulation of such information defines a natural informational
arrow of time.
Spacetime geometry responds not merely to energy but to the irreversible
narrowing of the set of physically possible futures.

The quantity $C$ denotes the physical complexity of the classical
record: the total surprisal $-\ln p(\text{history})$ of the complete
sequence of irreversible outcomes.
For locally uncorrelated R-events this reduces to the cumulative
surprisal
\begin{equation}
  C = \sum_i \bigl(-\ln p_i\bigr),
\end{equation}
where $p_i$ is the probability assigned by quantum mechanics to the
$i$-th R-process outcome prior to its occurrence.
This sum provides an operational lower bound on the full physical
complexity; equality holds when R-events are statistically independent,
which is justified by the definition of objective state reduction
\cite{Penrose1994}.
The full quantity coincides with the Kolmogorov complexity of the
classical record up to an additive constant bounded by the description
length of the laws of physics \cite{Kolmogorov1965}.
The derivation requires only the local rate $dC/d\tau$ under the
uncorrelated assumption; no computability is assumed, and the
lower-bound character of the sum does not affect the derivation.
Unitary quantum evolution contributes zero to $C$: each step of a
unitary evolution has $p_i = 1$, contributing zero surprisal.
Only irreversible transitions --- where $p_i < 1$ reflects genuine
quantum indeterminacy resolved into a definite outcome --- introduce
new terms to the sum.

The argument originates in a simple observation: $E = hf$ and $E = Mc^2$
share a left-hand side.
Setting them equal gives every isolated mass a characteristic frequency
$f = Mc^2/h$ --- the Compton frequency --- and raises a concrete
question: what is the mass doing at that frequency?
The answer proposed here is that it behaves as a Shannon source
\cite{Shannon1948}, producing irreversible classical information at that
rate.
The holographic principle then supplies the geometric link: information
requires physical expression as boundary area, and since it propagates
outward at $c$, area accumulates in the surrounding region.
The key physical insight is that information generated near the mass
spends time proportional to $r/c$ \emph{in transit} before reaching
a bounding surface at radius $r$.
More area accumulates per unit radius closer to the source, giving
gravity its $1/r$ character and producing the Schwarzschild radial
metric.

A note on terminology: throughout this paper we use \emph{isolated mass}
rather than \emph{particle}.
The argument applies equally to a proton, a neutron star, or the
observable universe.
The Compton scale $\lambda_C = \hbar/Mc$ is the relevant length scale
in each case, regardless of internal structure.

\smallskip
\noindent\textit{A note on the microscopic picture.}
The present paper treats mass as a rate of irreversible information
generation without specifying the microscopic mechanism responsible.
Companion papers develop a concrete model in which this generation
arises from the interaction of mass with vacuum information
threads --- topological objects whose density and crossing structure
constitute the local spacetime geometry.
The kinematic derivation here does not depend on that picture, but
the thread model provides a geometric interpretation of each step
and is offered as motivation for readers seeking a more concrete
substrate.

\subsection{Irreversible Processes and Information Generation}
\label{sec:irreversible}

Quantum dynamics contains two qualitatively different forms of
evolution \cite{Penrose1994}: reversible unitary evolution (U-process),
which preserves information content, and irreversible state reduction
(R-process), which produces definite classical outcomes and thereby
generates new classical information.

A simple example illustrates the principle.
Consider a 50/50 beam splitter illuminated by single photons at a rate
of $10^9$ per second.
Each photon is detected in either the transmitted or reflected arm,
producing a binary sequence of outcomes.
Because the events are statistically independent, this sequence contains
roughly $10^9$ bits of algorithmically incompressible classical
information after one second.
The mass equivalent implied by hypothesis~\eqref{eq:hypothesis} is
\begin{equation}
  M = \frac{\hbar}{2\pi c^2}\frac{dC}{d\tau}
    = \frac{1.055\times10^{-34}}{2\pi \times 9\times10^{16}}
      \times 10^9
    \approx 1.86\times10^{-43}\;\text{kg}.
\end{equation}
The geometric effect is real --- the surrounding area genuinely increases
by $4\ell_P^2$ per nat generated --- but undetectable at this rate.
The example illustrates the universal mechanism operating at a scale
far below the mass-vacuum interaction that dominates the curvature budget.
The framework asserts that matter continuously drives analogous processes
through its interaction with vacuum degrees of freedom at a rate set by
the Compton frequency.
The microscopic mechanism is not specified here and remains an open
question addressed in companion papers.

A potential confusion deserves direct address.
The complexity generation rate $dC/d\tau$ for one kilogram of matter is
vastly larger than the thermodynamic entropy of the same mass
($\sim 10^{25}$ nats for a kilogram of gas at standard conditions).
These are different ledgers.
Thermodynamic entropy counts the arrangement of matter's own microscopic
degrees of freedom.
The complexity generation rate counts the ongoing output of the
matter-vacuum interaction --- information flowing into the structure of
the surrounding vacuum, not into the internal state of the matter.
The large ratio is a consistency requirement rather than a liability:
the derivation produces the correct Schwarzschild geometry precisely
because the generation rate is set by mass-energy (via $\hbar$ and
$\alpha = 2\pi$), not by the thermodynamic state of the matter.
A rate comparable to thermodynamic entropy would produce no measurable
curvature.

\subsection{Relation to Prior Work}
\label{sec:prior}

The identification of quantum state reduction as the source of
irreversibility follows Penrose \cite{Penrose1994, HawkingPenrose1996},
who distinguished unitary from irreversible quantum evolution and
proposed that R-processes are connected to spacetime geometry.
Our framework inverts Penrose's causal direction: where Penrose proposed
that geometry \emph{drives} state reduction, we find that irreversible
transitions \emph{produce} geometry.
The R-process is primary; curvature is its accommodation.
We remain agnostic about the physical mechanism responsible for the
transition itself.

The holographic principle --- that the information content of a region
is bounded by its boundary area --- provides the bridge from information
generation to geometry.
Developed by 't~Hooft \cite{tHooft1993} and Susskind (see Other
References) and given precise form by Bekenstein's bound
\cite{Bekenstein1973, Bekenstein1981}, it is the second pillar of our
argument.
That entanglement entropy in quantum field theory scales with boundary
area rather than volume \cite{Srednicki1993} provides concrete evidence
that area-based information structure is a natural feature of quantum
spacetime, independently supporting assumption A2 below.

Several previous approaches have explored connections between gravity
and information.
The present framework differs from thermodynamic derivations of
gravitational dynamics in a key respect: rather than treating gravity
as an emergent entropic phenomenon, we identify mass itself as a rate
of irreversible information generation and interpret curvature as the
kinematic accommodation of that production.
The companion papers develop this distinction in detail.

\subsection{Hypothesis and Assumptions}
\label{sec:hypothesis}

The derivation rests on one central hypothesis and four assumptions.

\medskip
\noindent\textbf{H (Generation hypothesis).}
An isolated mass $M$ of pressureless dust generates irreversible
classical information at rate
$dC/d\tau = \alpha Mc^2/\hbar$,
where $\alpha$ is a single free parameter.
Section~2 establishes $\alpha = 2\pi$ from the Bekenstein bound alone,
without reference to the Schwarzschild metric.
The derivation of $g_{rr}$ in Section~3 then constitutes a prediction;
matching to the known Schwarzschild solution provides confirmation.

\medskip
\noindent\textbf{A1 (Propagation).}
Information propagates at the speed of light.
Classical information is encoded in physical states whose propagation
is bounded by $c$; R-process outputs cannot be transmitted
superluminally without violating relativistic causality.

\medskip
\noindent\textbf{A2 (Saturation).}
The vacuum near matter operates at holographic saturation: newly
generated information requires a corresponding increase in bounding
area.
A2 is an assumption of this framework; Srednicki \cite{Srednicki1993}
provides independent theoretical support from the area scaling of
entanglement entropy in quantum field theory, and Appendix~\ref{app:A2}
provides an independent quantitative motivation from an Olbers-style
integration of cosmological information flux.
We note explicitly that the successful derivation of the Schwarzschild
metric in Section~3 provides a posteriori consistency but not
independent confirmation of A2: a derivation that uses A2 cannot serve
as evidence for A2.
The strongest independent motivation for A2 is the cosmological
argument of Appendix~\ref{app:A2}, which gives holographic saturation
to within a factor of 3 under the Einstein--de Sitter approximation;
the exact calculation for the concordance cosmology is developed in
Paper~III.

\medskip
\noindent\textbf{A3 (Bekenstein bound).}
Each nat of information requires a minimum area of $4\ell_P^2$, where
$\ell_P = \sqrt{G\hbar/c^3}$.

\medskip
\noindent\textbf{A4 (Minimality).}
Nature produces no excess area beyond what the information requires.

\subsection{Departures from Standard Frameworks}
\label{sec:departures}

The central hypothesis has several consequences that depart from
standard physical assumptions.
We state them plainly here so that the reader can assess them
explicitly rather than encounter them as surprises.

\medskip
\noindent\textbf{Quantum evolution is not strictly unitary.}
R-processes are genuinely irreversible, not merely apparent
irreversibility arising from coarse-graining or loss of phase
information.
The framework treats state reduction as a primitive, not as something
to be derived from or reconciled with unitary evolution.
This is a deliberate departure from the Everettian programme; it
follows Penrose's insistence that R-processes are physically real,
though with the causal direction reversed.

A deeper geometric reason underlies this departure.
Unitary evolution preserves the inner product $\langle\psi|\phi\rangle$,
which presupposes a fixed measure and hence a fixed metric.
The Ricci tensor $R_{\mu\nu}$, by contrast, is the geometric record of
irreversible area-growth events and requires a dynamic metric.
These two requirements are mutually inconsistent.
By Gauss's \emph{Theorema Egregium}, a surface of nonzero Gaussian
curvature admits no isometric embedding in flat space: the surface cannot
be flattened without tearing.
Demanding that R-processes be reducible to U-processes is demanding
that a curved surface be globally flattenable --- possible in local
normal coordinates, impossible globally.
The ultraviolet divergences of perturbative quantum gravity are the
precise mathematical expression of this geometric incompatibility:
they are the infinities generated by the attempt to impose a flat
Hilbert-space inner product on a curved-spacetime dynamics.
This incompatibility has been established rigorously in the mathematical
physics literature: quantum field dynamics on generic curved spacetimes
is not unitarily implementable (see Related Works).
The algebraic approach to QFT was developed in part to accommodate
this failure by abandoning the Hilbert space apparatus.
The present framework takes a different lesson from the same result:
rather than preserving unitarity by modifying the formalism, we accept
that R-processes are genuinely non-unitary and identify curvature as
their geometric record.
The measurement problem is not a problem of interpretation.
It is a geometry problem.

\medskip
\noindent\textbf{Information is not conserved in the sense of a
fixed total.}
Each R-process produces a classical, objective outcome: a settled
fact added permanently to the universe's record.
In the sense of a museum conservator, information is never lost ---
settled outcomes are not undone --- but the collection is always
growing.
The total classical information content of the universe increases
monotonically.
This monotonic accumulation is the physical origin of the arrow of
time in the present framework.
Irreversibility deepens in stages: a measurement outcome is
irreversible in the practical sense (phase information gone);
information scrambled with many degrees of freedom is irreversible
in the thermodynamic sense; but the deepest form is causal ---
information crossing a horizon is placed permanently beyond the reach
of every future observer within the accessible universe.
Horizon-crossing is the universe's irrevocable commitment.

\medskip
\noindent\textbf{CPT symmetry is broken at the level of fundamental
dynamics.}
If R-processes are genuinely time-asymmetric, CPT violation is not
merely a thermodynamic artifact but a feature of the underlying
physics.
The expected magnitude, suppressed by approximately
$1/(\text{age of universe}) \approx 10^{-17}$--$10^{-18}$~s$^{-1}$,
corresponds to an energy scale of order $10^{-42}$~GeV --- roughly
24 orders of magnitude below the current best experimental bounds
from neutral kaon systems ($\sim 10^{-18}$~GeV; see Other References).
The prediction is therefore consistent with all current experimental
results and not accessible to near-term experiments.
Whether the suppression is uniform or concentrates near compact masses
where the local generation rate is high is a question for companion
papers and future experiment.

\medskip
\noindent\textbf{Spacetime curvature is kinematic rather than
dynamical.}
In standard general relativity, curvature is governed by field
equations: it is the dynamical response of spacetime to an
energy-momentum source.
In the present framework, curvature is a kinematic consequence of
information-theoretic constraints --- the geometric shape that
information generation necessarily imposes on the surrounding
spacetime, in the same way that Lorentz contraction is a kinematic
consequence of the invariance of $c$ rather than a dynamical force.
Whether this reframing has observable consequences beyond reproducing
the standard results is a question the companion papers begin to
address.

\subsection{Scope of the Present Paper}
\label{sec:scope}

The aim of the present paper is to derive the exterior metric for
spherically symmetric mass distributions from H and A1--A4, without
invoking Einstein's field equations.
This is carried out for the uncharged case (Section~3, recovering
Schwarzschild) and for the charged case (Section~4, recovering
Reissner--Nordstr\"om).
An argument for the temporal component $g_{tt}$ is presented as a
conjecture open to further scrutiny.
Questions concerning the microscopic origin of the information
generation process, the derivation of the Einstein equations, and
cosmological information balance are addressed in companion papers.

%% ============================================================
\section{Fixing the Free Parameter: $\alpha = 2\pi$ from the
Bekenstein Bound}
\label{sec:alpha}
%% ============================================================

The hypothesis H contains one free parameter $\alpha$.
We show that this parameter is fixed uniquely by the Bekenstein bound,
without invoking the Schwarzschild metric.
The derivation of $g_{rr}$ in Section~3 is therefore a prediction, not
a fitting exercise; the Schwarzschild matching confirms it.

\subsection{The Bekenstein Bound}

The Bekenstein bound limits the information content of a region of
radius $r$ containing energy $E = Mc^2$.
In natural units ($\hbar = c = 1$):
\begin{equation}
  K_{\max} = 2\pi Mr.
  \label{eq:bekenstein-bound}
\end{equation}
The factor $2\pi$ appears in Bekenstein's original derivation
\cite{Bekenstein1981}, reflecting the relationship between energy,
radius, and holographic capacity.

\subsection{Information in Transit with Free Parameter}

By H and A1, the information currently in transit between the source and
a sphere at radius $r$ (in natural units $\hbar = c = 1$) is, for
general $\alpha$:
\begin{equation}
  \Delta C = \frac{dC}{d\tau}\cdot\frac{r}{c} = \alpha Mr.
  \label{eq:transit-natural}
\end{equation}
Here $\alpha$ is the free parameter of H, not yet fixed.

\subsection{Fixing $\alpha$ by Saturation}

We require that the information in transit saturate the Bekenstein
bound at every radius:
\begin{equation}
  \Delta C = K_{\max}.
  \label{eq:saturation-condition}
\end{equation}
Substituting equations~\eqref{eq:bekenstein-bound}
and~\eqref{eq:transit-natural}:
\begin{equation}
  \alpha Mr = 2\pi Mr \quad \text{for all } r,
\end{equation}
which requires
\begin{equation}
  \alpha = 2\pi.
  \label{eq:alpha-fixed}
\end{equation}
This fixes the free parameter of hypothesis~H from the Bekenstein
bound alone, without reference to the Schwarzschild metric.
The exact factor of $2\pi$ is not a fitted value; it is the only
value consistent with holographic saturation at every radius.

With $\alpha = 2\pi$ established, the information generation rate is:
\begin{equation}
  \frac{dC}{d\tau} = \frac{2\pi Mc^2}{\hbar},
  \label{eq:rate-fixed}
\end{equation}
and the information in transit to a sphere at radius $r$ is:
\begin{equation}
  \Delta C = \frac{2\pi Mcr}{\hbar}.
  \label{eq:transit-fixed}
\end{equation}

\subsection{Summary and Forward-Reference}

The requirement that information in transit saturates the Bekenstein
bound at every radius fixes $\alpha = 2\pi$ uniquely, with no reference
to the Schwarzschild metric.
With $\alpha = 2\pi$, the rate $dC/d\tau = 2\pi Mc^2/\hbar$:
reproduces the Schwarzschild radial metric (Section~3.3) as a
prediction; saturates the Bekenstein bound at every radius (this
section); and yields the factor $2\pi$ as an inevitable consistency
requirement rather than a free choice.

\medskip
\noindent\textbf{Saturation as a governing principle.}
Mass generates information at exactly the maximum rate compatible with
the geometric capacity of its surroundings.
This ``just-in-time'' character ensures consistency without fine-tuning.

\medskip
\noindent\textbf{Compton scale validation.}
The equality holds at all radii, including the Compton wavelength
$\lambda_C = \hbar/Mc$, providing a natural consistency check at the
quantum scale.

%% ============================================================
\section{Derivation of the Exterior Schwarzschild Metric}
\label{sec:schwarzschild}
%% ============================================================

We derive the exterior Schwarzschild metric from hypothesis H and
assumptions A1--A4, using $\alpha = 2\pi$ as established in Section~2.
No field equations, curvature tensors, or prior knowledge of the
metric's form are required.

\subsection{Area Response to Information Growth}
\label{sec:area-response}

The derivation rests on the following chain of reasoning.
Classical information is generated through irreversible quantum state
reductions.
Because such reductions are non-unitary, classical information
increases monotonically, providing a natural informational arrow of
time.
Information contained within a spatial region is subject to the
Bekenstein--Hawking bound: the maximum entropy associated with a
region is proportional to the area of its bounding surface,
\begin{equation}
  S_{\max} = \frac{A}{4\ell_P^2}.
  \label{eq:bekenstein-hawking}
\end{equation}
If classical information within a region increases through irreversible
state reductions, the holographic bound implies that the bounding area
must increase accordingly.
This requirement links information growth directly to geometry.

A preliminary calculation shows that the ratio of holographic
information to holographic capacity yields the exact Schwarzschild
parameter.
The Bekenstein bound for mass $M$ within radius $r$ (natural units
$\hbar = c = 1$) is
\begin{equation*}
  S_{\rm Bekenstein} = 2\pi Mr.
\end{equation*}
The holographic capacity of the bounding sphere of radius $r$
(with $\ell_P^2 = G$ in natural units) is
\begin{equation*}
  S_{\max} = \frac{4\pi r^2}{4G} = \frac{\pi r^2}{G}.
\end{equation*}
The fractional saturation is therefore, exactly,
\begin{equation*}
  \frac{S_{\rm Bekenstein}}{S_{\max}}
  = \frac{2\pi Mr}{\pi r^2/G}
  = \frac{2GM}{r}
  = \frac{r_S}{r},
\end{equation*}
where $r_S = 2GM/c^2$ is the Schwarzschild radius.
This is not a scaling result: the numerical prefactors work out
exactly, with no residual dimensionless factor.
The gravitational field strength at radius $r$ is precisely the
degree to which the region's informational content saturates its
holographic capacity.
It is exactly this ratio $r_S/r$ that controls every deviation from
flat space in both components of the Schwarzschild metric:
$g_{tt} = -(1 - r_S/r)$ and $g_{rr} = (1 - r_S/r)^{-1}$.

For a general static spherically symmetric line element
\cite{Wald1984}
\begin{equation}
  ds^2 = -A(r)\,dt^2 + B(r)\,dr^2 + r^2\,d\Omega^2,
  \label{eq:general-metric}
\end{equation}
the radial metric coefficient $B(r)$ controls the physical length
associated with a coordinate radial displacement.
The derivation in Sections~3.3 and~3.4 makes the informational
origin of $B(r)$ precise.

\subsection{Setup and Coordinate Conventions}
\label{sec:setup}

Consider a spherically symmetric, stationary mass $M$ of pressureless
dust.
We work in the exterior vacuum region $r > R$, where $R$ is the radius
of the matter distribution, and the geometry approaches flat Minkowski
spacetime as $r \to \infty$.
We adopt Schwarzschild coordinates $(t, r, \theta, \phi)$ \cite{Wald1984}.
The Schwarzschild radius is
\begin{equation}
  r_S = \frac{2GM}{c^2}.
  \label{eq:rs}
\end{equation}
This is a geometric length scale; it does not imply the object is a
black hole.
We proceed in two steps: $g_{rr}$ from area excess (Section~3.3),
then the argument for $g_{tt}$ (Section~3.4).

\subsection{The Radial Component: Area Excess from Information
in Transit}
\label{sec:grr}

With $\alpha = 2\pi$ established in Section~2, hypothesis H gives the
information generation rate
\begin{equation}
  \frac{dC}{d\tau} = \frac{2\pi Mc^2}{\hbar}.
  \label{eq:gen-rate}
\end{equation}
By A1, information propagates outward at $c$.
In the stationary configuration, the information currently in transit
between the source and a sphere at coordinate radius $r$ was generated
during the preceding proper time interval $r/c$, giving
\begin{equation}
  \Delta C = \frac{dC}{d\tau}\cdot\frac{r}{c}
           = \frac{2\pi Mcr}{\hbar}.
  \label{eq:transit}
\end{equation}
By A2 (saturation), this information requires geometric expression as
area.
By A3 (Bekenstein), each nat requires $4\ell_P^2$ of area.
By A4 (minimality), the area excess is exactly
\begin{equation}
  \Delta A = 4\ell_P^2\,\Delta C
           = 4\frac{G\hbar}{c^3}\cdot\frac{2\pi Mcr}{\hbar}
           = \frac{8\pi GMr}{c^2}
           = 2 r_S r.
  \label{eq:deltaA}
\end{equation}
We define the radial coordinate $r$ as the areal coordinate: the
physical area of a sphere at coordinate $r$ is $A_{\rm phys} = 4\pi r^2$.
The area that sphere would have absent the mass's information
contribution is
\begin{equation}
  A_{\rm unstretched} = 4\pi r^2 - 2 r_S r.
  \label{eq:unstretched}
\end{equation}
The radial metric component is the ratio of physical to unstretched
area:
\begin{equation}
  g_{rr} = \frac{4\pi r^2}{4\pi r^2 - 2 r_S r}
          = \frac{1}{1 - r_S/r}.
  \label{eq:grr}
\end{equation}
This is the radial component of the Schwarzschild metric, derived from
information accounting alone.
Each step follows from H (with $\alpha = 2\pi$ from Section~2) and
A1--A4 in sequence.
Matching to the known Schwarzschild metric $g_{rr} = (1 - r_S/r)^{-1}$
confirms the result as a prediction.

\medskip
\noindent\textbf{Proper radial length and holographic saturation.}
The ratio of coordinate displacement to proper radial length is
\begin{equation*}
  \frac{dr}{d\ell}
  = \frac{1}{\sqrt{g_{rr}}}
  = \sqrt{1 - \frac{r_S}{r}}
  = \sqrt{1 - \frac{S_{\rm Bekenstein}}{S_{\max}}}.
\end{equation*}
The ratio of coordinate to proper radial length is the square root
of the \emph{unsaturated} holographic fraction.
At $r = r_S$, saturation is complete: $dr/d\ell \to 0$, and an
infinite proper distance separates any exterior point from the horizon.
The horizon is the surface of complete holographic saturation, and the
divergence of $g_{rr}$ there is not a coordinate artifact but a direct
geometric expression of that saturation.

\subsection{The Temporal Component: A Conjecture and an Invitation}
\label{sec:gtt}

The derivation of $g_{rr}$ is kinematic: it follows from counting
information in transit.
The temporal component $g_{tt}$ requires a further argument.
We present the strongest argument we have, identify the step that
remains a conjecture, and invite community scrutiny.
An independent derivation of $g_{tt}$ via the Einstein equations
appears in Paper~II.

The configuration is stationary, so the spacetime admits a timelike
Killing vector $\partial/\partial t$.
The energy of any null ray --- the carrier of information --- is
conserved along each trajectory in the exterior:
\begin{equation}
  E = -g_{tt}\frac{dt}{d\lambda} = \text{const along each null ray.}
  \label{eq:killing-energy}
\end{equation}
In the source-free exterior, no information is generated.
The integrated information flux through any concentric sphere is
therefore constant:
\begin{equation}
  \Phi(r) = \text{const} \qquad \text{for all } r > R.
  \label{eq:flux-conservation}
\end{equation}
The most general static spherically symmetric metric takes the form
\cite{Wald1984}
\begin{equation}
  ds^2 = -A(r)\,dt^2 + B(r)\,dr^2 + r^2\,d\Omega^2,
  \label{eq:metric-general}
\end{equation}
with $B(r) = (1 - r_S/r)^{-1}$ already determined by Section~3.3.
We write
\begin{equation}
  A(r) = e^{2\Psi(r)}\left(1 - \frac{r_S}{r}\right)
  \label{eq:Ar}
\end{equation}
for a free function $\Psi(r)$.
Asymptotic flatness requires $\Psi(r) \to 0$ as $r \to \infty$;
the general expansion is $\Psi(r) = \beta/r + \gamma/r^2 + \cdots$
Each nonzero coefficient introduces a free parameter that must be
sourced by some physical field.
In the present setup --- vacuum exterior, single length scale $r_S$
already fixed, no additional fields --- no such source is present.

\medskip
\noindent\textbf{Conjecture:} $\Psi(r) = 0$ identically.

\medskip
By A4 (minimality), no geometric structure is produced beyond what
the information constraints require.
A nonzero $\Psi(r)$ would modify the proper time measured by stationary
clocks at radius $r$, changing the rate at which information generated
near the source is red- or blue-shifted as it propagates outward.
This modified propagation rate would require additional area
accommodation beyond what the information in transit already demands ---
area sourced by nothing in the physical setup, violating A4.
Therefore $g_{tt} = -(1 - r_S/r)$.

\medskip
\noindent\textbf{Where the argument needs scrutiny.}
A4 is doing substantial work here.
In GR, $\Psi = 0$ follows from the vacuum Einstein equations via
Birkhoff's theorem --- a rigorous result.
We reach the same conclusion by a different route.
The open problem is to show rigorously that A4 excludes all unsourced
free parameters in the temporal component, without implicitly assuming
the conclusion.
We present this as a conjecture; the companion paper (Paper~II) provides
an independent derivation that may serve as confirmation pending a
direct proof.

\subsection{The Exterior Schwarzschild Metric}
\label{sec:schwarzschild-metric}

If the conjecture of Section~3.4 holds, the complete line element in
the exterior region $r > R$ is
\begin{equation}
  ds^2 = -\!\left(1 - \frac{r_S}{r}\right)c^2\,dt^2
         + \left(1 - \frac{r_S}{r}\right)^{-1}dr^2
         + r^2\!\left(d\theta^2 + \sin^2\!\theta\,d\phi^2\right),
  \label{eq:schwarzschild}
\end{equation}
with $r_S = 2GM/c^2$ and $\alpha = 2\pi$.
The radial component is exact.
The temporal component depends on the conjecture of Section~3.4.
Together they constitute the exterior Schwarzschild metric.

%% ============================================================
\section{Extension to Charged Matter:
Derivation of the Exterior Metric}
\label{sec:charged}
%% ============================================================

We now apply the framework to a spherically symmetric, stationary mass
carrying total electric charge $Q$.
The aim is to derive the exterior metric using H and A1--A4 alone,
with no additional assumptions.
In particular, Maxwell's equations, the Faraday tensor $F_{\mu\nu}$,
and the electromagnetic stress-energy tensor $T_{\mu\nu}^{\rm EM}$ play
no role in what follows.

\subsection{The Enclosed Energy at Radius $r$}
\label{sec:rn-setup}

The uncharged derivation took the information-generating energy to be
$Mc^2$, independent of $r$.
For a charged body this requires care, because the electromagnetic
field carries energy distributed throughout the exterior of the source.

The total ADM mass $M$ --- the conserved energy measured at spatial
infinity --- includes both the irreducible (rest) mass $M_{\rm irr}$
and the total electromagnetic self-energy of the field:%
\footnote{For weakly charged astrophysical objects,
$E_{\rm field,total}/c^2 \ll M_{\rm irr}$ and $M \approx M_{\rm irr}$
to good approximation.
For near-extremal objects the distinction is physically significant.
The electron, with $Q_e/m_e \sim 10^{21}$ in Planck units, has
electromagnetic self-energy that dominates its rest energy in classical
terms; the regularisation of this divergence is an open problem for
the framework and for classical electrodynamics alike.
The present derivation assumes the decomposition
\eqref{eq:adm-decomp} is well-defined, which it is for any finite
charge distribution with a characteristic radius.}
\begin{equation}
  M = M_{\rm irr} + \frac{E_{\rm field,total}}{c^2}.
  \label{eq:adm-decomp}
\end{equation}
The electromagnetic energy density of a Coulomb field at radius $r'$ is
\begin{equation}
  u(r') = \frac{k_e Q^2}{8\pi r'^4},
  \label{eq:em-density}
\end{equation}
where $k_e = 1/(4\pi\varepsilon_0)$ in SI units.

The field energy \emph{exterior} to radius $r$ is obtained by
integrating inward from infinity --- where the field energy is zero
and where the ADM mass is defined --- to $r$:
\begin{equation}
  E_{\rm field,ext}(r)
    = \int_r^\infty u(r')\,4\pi r'^2\,dr'
    = \frac{k_e Q^2}{2}\int_r^\infty r'^{-2}\,dr'
    = \frac{k_e Q^2}{2r}.
  \label{eq:field-ext}
\end{equation}
This integral is convergent at infinity and well-defined for all
$r > 0$.
We integrate from $r = \infty$ inward because this is the natural
direction: the ADM mass is the boundary value, and we subtract field
energy as we move inward past each shell.

The energy \emph{enclosed} within radius $r$ --- the portion of the
total mass-energy that has already arrived at or inside the sphere ---
is therefore
\begin{equation}
  E_{\rm enc}(r) = Mc^2 - \frac{k_e Q^2}{2r}.
  \label{eq:e-enc}
\end{equation}
This is the quantity that drives information generation as seen by the
bounding sphere at $r$.

\subsection{Information from the Enclosed Energy}
\label{sec:rn-info}

By H with $\alpha = 2\pi$ (Section~2), the information generation rate
associated with the enclosed energy is
\begin{equation}
  \frac{dC}{d\tau}
    = \frac{2\pi}{\hbar}\,E_{\rm enc}(r)
    = \frac{2\pi}{\hbar}\!\left(Mc^2 - \frac{k_e Q^2}{2r}\right).
  \label{eq:rn-rate}
\end{equation}
By A1, this information accumulates over the light-crossing time $r/c$:
\begin{align}
  \Delta C
    &= \frac{dC}{d\tau}\cdot\frac{r}{c} \notag \\
    &= \frac{2\pi}{\hbar c}\!\left(Mc^2 r - \frac{k_e Q^2}{2}\right) \notag \\
    &= \frac{2\pi}{\hbar}\!\left(Mcr - \frac{k_e Q^2}{2c}\right).
  \label{eq:rn-transit}
\end{align}
The $r$-dependence in $E_{\rm enc}$ --- the $1/r$ field energy term ---
combines with the factor of $r$ from the light-crossing time to produce
a term $k_e Q^2/(2c)$ that is \emph{independent of $r$}.
This cancellation is the key structural fact of the charged derivation:
the charge contribution to the accumulated information is constant in
$r$, which, when spread over the sphere of area $4\pi r^2$, generates
a $1/r^2$ metric term.

\subsection{Area Excess and Metric Coefficients}
\label{sec:rn-area}

By A2, A3, and A4, the area required to encode $\Delta C$ is
\begin{align}
  \Delta A
    &= 4\ell_P^2\,\Delta C \notag \\
    &= 4\frac{G\hbar}{c^3}\cdot\frac{2\pi}{\hbar}
       \!\left(Mcr - \frac{k_e Q^2}{2c}\right) \notag \\
    &= \frac{8\pi G}{c^2}\!\left(Mr - \frac{k_e Q^2}{2c^2}\right) \notag \\
    &= \frac{8\pi GMr}{c^2} - \frac{4\pi G k_e Q^2}{c^4}.
  \label{eq:rn-area}
\end{align}
Distributing this area excess over the sphere of area $4\pi r^2$:
\begin{equation}
  \frac{\Delta A}{4\pi r^2}
    = \frac{2GM}{rc^2} - \frac{G k_e Q^2}{c^4 r^2}.
  \label{eq:rn-ratio}
\end{equation}
By the same area-stretching argument as in Section~3.3, the radial
metric component is
\begin{equation}
  g_{rr} = \frac{1}{1 - 2GM/(rc^2) + G k_e Q^2/(c^4 r^2)}.
  \label{eq:rn-grr}
\end{equation}
By the same flux-conservation conjecture as in Section~3.4,
$g_{tt}\,g_{rr} = -1$ in the exterior vacuum region, giving
\begin{equation}
  g_{tt} = -\!\left(1 - \frac{2GM}{rc^2} + \frac{G k_e Q^2}{c^4 r^2}\right).
  \label{eq:rn-gtt}
\end{equation}

\subsection{The Exterior Charged Metric}
\label{sec:rn-metric}

Introducing the charge length scale
\begin{equation}
  r_Q^2 \equiv \frac{G k_e Q^2}{c^4}
             = \frac{GQ^2}{4\pi\varepsilon_0 c^4}
  \label{eq:rq}
\end{equation}
alongside the Schwarzschild radius $r_S = 2GM/c^2$, the complete line
element in the exterior region $r > R$ is
\begin{equation}
  ds^2 = -f(r)\,c^2\,dt^2
         + f(r)^{-1}\,dr^2
         + r^2\!\left(d\theta^2 + \sin^2\!\theta\,d\phi^2\right),
  \label{eq:charged-metric}
\end{equation}
where
\begin{equation}
  f(r) = 1 - \frac{r_S}{r} + \frac{r_Q^2}{r^2}.
  \label{eq:rn-f}
\end{equation}
For $Q = 0$, $r_Q = 0$ and equation~\eqref{eq:charged-metric} reduces
to the Schwarzschild metric~\eqref{eq:schwarzschild}, as required.

This metric was derived without invoking Maxwell's equations, the
Faraday tensor, or the electromagnetic stress-energy tensor.
The $r_Q^2/r^2$ correction arises entirely from the geometry of energy
enclosure: the Coulomb field energy exterior to $r$ is subtracted from
$M$ before computing the information budget, and the resulting
$r$-dependent correction to the enclosed energy generates a constant
area term that spreads as $1/r^2$ across the sphere.
The sign of this correction is positive (repulsive), opposing the
gravitational attraction of the mass term, consistent with the
electrostatic self-repulsion of a charged distribution.

\subsection{Comparison with the Reissner--Nordstr\"om Solution}
\label{sec:rn-comparison}

Equation~\eqref{eq:charged-metric} is the
\emph{Reissner--Nordstr\"om metric} \cite{Wald1984} for a charged,
non-rotating, spherically symmetric body.
In standard general relativity it is obtained by solving the
Einstein--Maxwell equations $G_{\mu\nu} = (8\pi G/c^4)T_{\mu\nu}$
with the electromagnetic stress-energy tensor $T_{\mu\nu}^{\rm EM}$,
a calculation that requires knowing the structure of
$F_{\mu\nu}F^{\mu\nu}$ and tracing through coupled partial
differential equations.

The present derivation requires none of this.
The same two-parameter framework --- one hypothesis H, four assumptions
A1--A4, no additional input --- that produces the Schwarzschild metric
for $Q = 0$ produces the Reissner--Nordstr\"om metric for $Q \neq 0$.
The transition from Schwarzschild to Reissner--Nordstr\"om is not a
new assumption but a consequence of correctly accounting for how charge
distributes its mass-energy in space.

This structural unity provides evidence against an accidental-concordance
reading of the Schwarzschild derivation of Section~3.
A numerical coincidence would not extend cleanly to the two-parameter
case without modification.
The correct $Q^2/r^2$ coefficient, correct sign, and correct recovery
of Schwarzschild in the $Q \to 0$ limit all follow from the same
physical picture.

\subsection{Horizons and the Extremal Limit}
\label{sec:rn-horizons}

Setting $f(r) = 0$ gives the horizon condition.
In terms of $r_S$ and $r_Q$, the two solutions are
\begin{equation}
  r_\pm = \frac{r_S}{2}\!\left[1 \pm
          \sqrt{1 - \left(\frac{2r_Q}{r_S}\right)^{\!2}}\,\right].
  \label{eq:rn-horizons}
\end{equation}
For $2r_Q < r_S$ (sub-extremal), there are two horizons: the outer
event horizon $r_+$ and the inner Cauchy horizon $r_-$.
For $2r_Q = r_S$ (extremal), the two horizons coincide at
$r = r_S/2 = GM/c^2$; in this limit $M_{\rm irr} \to 0$ and the ADM
mass is dominated by electromagnetic self-energy, representing a
purely electromagnetic object.
For $2r_Q > r_S$ (super-extremal), there are no horizons; this would
require negative irreducible mass, which is apparently not realized
in nature.

%% ============================================================
\section{Discussion and Conclusions}
%% ============================================================

\subsection{Summary of Results}

We have derived the exterior metric for spherically symmetric mass
distributions from one hypothesis (H) and four assumptions (A1--A4),
without invoking Einstein's field equations.

The free parameter $\alpha$ in H is fixed to $\alpha = 2\pi$ by the
Bekenstein bound alone (Section~2), without reference to any known
metric.
With $\alpha$ so fixed:

\begin{itemize}

\item The radial component of the exterior Schwarzschild metric,
$g_{rr} = (1 - r_S/r)^{-1}$, follows as a prediction (Section~3.3),
confirmed by matching to the known solution.
The derivation is exact within the stated assumptions.

\item The temporal component $g_{tt} = -(1 - r_S/r)$ follows from a
well-motivated conjecture (Section~3.4); an independent derivation
appears in Paper~II.

\item The exterior metric for a spherically symmetric charged mass
(Section~4) is recovered exactly, including the correct $Q^2/r^2$
coefficient and sign, without invoking Maxwell's equations or the
electromagnetic stress-energy tensor.

\end{itemize}

The structural unity of the uncharged and charged derivations within
a single framework, with no additional assumptions introduced for the
charged case, is a non-trivial consistency check on the framework.

\subsection{Extensions in Companion Papers}

The following topics are developed in companion papers.

\medskip
\noindent\textbf{Einstein equations (Paper~II).}
The full Einstein field equations for vacuum and pressureless dust,
derived from Screen Theory premises.
$g_{tt}$ is recovered here as a theorem rather than a conjecture.

\medskip
\noindent\textbf{Cosmology (Paper~III).}
Global information balance, Machian origin of mass-energy, and the
cosmological constant reinterpreted in information-theoretic terms.

\medskip
\noindent\textbf{Horizons and black holes (Paper~IV).}
Behavior at and inside $r = r_S$, information export rates, and the
Page curve in information-theoretic terms.

\subsection{Conceptual Implications}
\label{sec:implications}

The derivation suggests a fundamental reinterpretation.
Mass is not a source of curvature in the sense of a cause producing an
effect; it \emph{is} a rate of information generation, and curvature
is the geometric shape that rate imposes on its surroundings.
This makes precise Wheeler's \cite{Wheeler1990} ``it from bit'' vision:
geometry (\emph{that}) from information generation (\emph{nat}).

The causal order is inverted relative to standard GR\@.
In GR, mass-energy is a source term in dynamical field equations and
curvature is the solution.
Here, irreversible information generation is primary and geometry is
its kinematic accommodation.
The Compton frequency $Mc^2/\hbar$, typically interpreted as a quantum
oscillation frequency, acquires new significance as the characteristic
rate of irreversible constraint production.

\medskip
\noindent\textbf{Newton's gravitational constant reinterpreted.}
From the derivation chain, $G$ can be written as
\begin{equation*}
  G = \frac{\ell_P^2 c^3}{\hbar},
\end{equation*}
where $4\ell_P^2$ converts area to information (A3) and $\hbar$
converts action to information.
$G$ is therefore not a fundamental coupling constant but an
\emph{exchange rate} between Planck area and quantum of action.
The factor $c^3$ adjusts units only.
$G$'s apparent smallness is not a mystery requiring explanation but
a consequence of the ratio between two independently motivated
fundamental quantities.

\medskip
\noindent\textbf{Charge as a redistribution of enclosed energy.}
The Reissner--Nordstr\"om result of Section~4 has an implication
for how we understand the role of charge in gravity.
Charge does not add a new source term to the curvature; it redistributes
the energy that is already counted in the ADM mass, spreading some of
it into the exterior field.
The curvature at radius $r$ responds to the energy that has arrived
inside $r$, not to the total.
In this picture, the $Q^2/r^2$ term is not an electromagnetic
correction to gravity but a consequence of the radial profile of energy
enclosure.

\subsection{An Open Problem and an Invitation}

The derivation of $g_{rr}$ is rigorous within the stated assumptions.
The argument for $g_{tt}$ rests on a step --- minimality excluding
unsourced free parameters --- that we present as a conjecture.
A rigorous proof of this step, a counterexample forcing modification
of the framework, or an alternative derivation closing the gap, would
each be a valuable contribution.
The companion paper (Paper~II) provides an independent route to $g_{tt}$
via the Einstein equations, which may serve as confirmation of the
conjecture pending a direct proof.

\subsection{Forward-Looking Conjectures}

\noindent\textbf{Microscopic thread structure.}
The present paper derives the Schwarzschild and Reissner--Nordstr\"om
geometries from information accounting without specifying the
microscopic objects that carry and generate the information.
Companion papers develop a model in which irreversible information
generation arises from the interaction of mass with vacuum information
threads --- topological objects whose crossing density and biographical
structure constitute local spacetime geometry.
In that picture, mass is a persistent pattern of thread crossings at
the Compton rate, and curvature is the area accumulated by threads
propagating outward from those crossings.
The thread model provides a concrete realization of hypothesis~H and
a geometric interpretation of each step in the derivation above.

\medskip
\noindent\textbf{Rotating matter: the Kerr metric.}
The natural next extension is to rotating, uncharged matter.
The Kerr metric introduces a second parameter (angular momentum $J$)
alongside $M$.
Whether the framework reproduces the Kerr solution from information
accounting alone --- without the Boyer--Lindquist ansatz or Penrose's
$r$-complex trick --- is an open problem.

\medskip
\noindent\textbf{Shape independence.}
Bekenstein's bound assumes nothing about interior mass distribution.
The present argument may inherit this insensitivity; the spherically
symmetric case may be the simplest instance of something more general.

\medskip
\noindent\textbf{Screen overlap and Newton's law.}
Newtonian gravity and its relativistic extension may admit a unified
interpretation via holographic screen overlap between masses, with the
nonlinearity of Einstein's equations corresponding to the breakdown of
the small-overlap approximation.

\medskip
\noindent\textbf{Cosmological information accumulation.}
The accumulation of classical information on cosmic horizons --- the
irreversible deposit of settled outcomes beyond the reach of all future
observers --- may drive large-scale spacetime dynamics.
This connection, including the Machian origin of mass-energy and the
derivation of $E = Mc^2$ from Screen Theory premises, is developed in
Paper~III.

%% ============================================================
\section*{Acknowledgments}
%% ============================================================

The author thanks Joel Welling for illuminating early conversations
about the holographic principle, Bob Carlitz for valuable feedback,
and Finula McCaul for valuable conversations.

%% ============================================================
\section*{Citations}
%% ============================================================

\begin{thebibliography}{99}

\bibitem{Bekenstein1973}
        Bekenstein, J.\,D.\ (1973).
        Black holes and entropy.
        \textit{Physical Review D}, 7(8), 2333--2346.

\bibitem{Bekenstein1981}
        Bekenstein, J.\,D.\ (1981).
        Universal upper bound on the entropy-to-energy ratio for
        bounded systems.
        \textit{Physical Review D}, 23(2), 287--298.

\bibitem{HawkingPenrose1996}
        Hawking, S.\,W.\ and Penrose, R.\ (1996).
        \textit{The Nature of Space and Time}.
        Princeton University Press.

\bibitem{Kolmogorov1965}
        Kolmogorov, A.\,N.\ (1965).
        Three approaches to the quantitative definition of information.
        \textit{Problems of Information Transmission}, 1(1), 1--7.

\bibitem{Mach1883}
        Mach, E.\ (1883).
        \textit{Die Mechanik in ihrer Entwicklung historisch-kritisch
        dargestellt}.
        Brockhaus, Leipzig.
        [English translation: \textit{The Science of Mechanics},
        Open Court, 1893.]

\bibitem{Penrose1994}
        Penrose, R.\ (1994).
        \textit{Shadows of the Mind}.
        Oxford University Press.

\bibitem{Shannon1948}
        Shannon, C.\,E.\ (1948).
        A mathematical theory of communication.
        \textit{Bell System Technical Journal},
        27(3), 379--423 and 27(4), 623--656.

\bibitem{Srednicki1993}
        Srednicki, M.\ (1993).
        Entropy and area.
        \textit{Physical Review Letters}, 71(5), 666--669.

\bibitem{tHooft1993}
        't~Hooft, G.\ (1993).
        Dimensional reduction in quantum gravity.
        arXiv:gr-qc/9310026.

\bibitem{Wald1984}
        Wald, R.\,M.\ (1984).
        \textit{General Relativity}.
        University of Chicago Press.

\bibitem{Wheeler1990}
        Wheeler, J.\,A.\ (1990).
        Information, physics, quantum: The search for links.
        In Zurek, W.\,H.\ (ed.),
        \textit{Complexity, Entropy, and the Physics of Information}.
        Addison-Wesley.

\end{thebibliography}

%% ============================================================
\section*{Other References}
%% ============================================================

\noindent\textit{A note on methodology.}
The following works were identified as closely related to the present
paper by an AI assistant (Claude, Anthropic, March 2026), on the basis
of semantic content matching against the physics literature, or are
works the authors are aware of but have not yet read.
The authors cannot vouch for the accuracy of the AI's assessments
or for the correctness of specific claims attributed to these works.

We list them here because the standard convention of citing only works
one has actually read creates a systematic gap in the scholarly record.
An AI can identify semantic neighbors in the literature rapidly and
reliably; human authors are better placed to evaluate whether a work
is correct and relevant in depth.
Separating these two tasks --- neighbor identification from depth
evaluation --- serves science better than conflating them under a
single citation convention.

We propose that \textbf{Other References} sections, explicitly
identified as AI-assisted or colleague-assisted rather than
author-read, constitute a distinct and honest category of scholarly
attribution.
Such sections provide readers with a navigable neighborhood of the
paper's ideas without misrepresenting the authors' familiarity with
that neighborhood.
We invite the community to adopt this convention.

\smallskip
\noindent\textit{Readers who find errors in the AI assessments, or who
are aware of additional closely related works, are encouraged to
contact the authors.}

\bigskip

\begin{enumerate}

  \item Arageorgis, A., Earman, J., and Ruetsche, L.\ (2002).
        Weyling the time away: the non-unitary implementability of
        quantum field dynamics on curved spacetime.
        \textit{Studies in the History and Philosophy of Modern Physics},
        33(2), 151--184.
        \textit{[AI assessment: establishes rigorously that QFT dynamics
        on generic curved spacetimes is not unitarily implementable ---
        the dilemma identified here is directly relevant to the
        departures from standard frameworks discussed in Section~1.5.
        The present framework's response to this dilemma differs from
        the algebraic approach advocated in that paper.]}

  \item Verlinde, E.\ (2011).
        On the origin of gravity and the laws of Newton.
        \textit{Journal of High Energy Physics}, 2011(4), 29.
        \textit{[AI assessment: derives gravitational dynamics from
        entropic considerations, sharing the broad programme of
        connecting gravity to information. The present framework
        differs in identifying mass (not gravity) as the fundamental
        information-theoretic quantity, and in treating curvature as
        kinematic rather than entropic. The precise distinction is
        developed in the companion papers.]}

  \item Jacobson, T.\ (1995).
        Thermodynamics of spacetime: the Einstein equation of state.
        \textit{Physical Review Letters}, 75(7), 1260--1263.
        \textit{[AI assessment: derives the Einstein equations from
        thermodynamic considerations applied to local Rindler horizons.
        Related in spirit to the present derivation but different in
        both starting point and method: Jacobson assumes the Clausius
        relation as fundamental, whereas the present framework derives
        geometry from information generation rates.]}

  \item Bilson-Thompson, S.\,O.\ (2005).
        A topological model of composite preons.
        arXiv:hep-ph/0503213.
        \textit{[AI assessment: represents Standard Model first-generation
        fermions and gauge bosons as specific braids of three ribbons ---
        directly relevant to the thread/braid picture developed in the
        companion papers and flagged as an open correspondence problem
        in the present programme.]}

  \item Penrose, R.\ (1971).
        Angular momentum: an approach to combinatorial space-time.
        In Bastin, T.\ (ed.),
        \textit{Quantum Theory and Beyond}.
        Cambridge University Press.
        \textit{[AI assessment: the original twistor/spin-network paper
        proposing that spacetime geometry should emerge from discrete
        combinatorial structure --- the intellectual ancestor of
        approaches (including the present one) that treat spacetime
        points as derived rather than primitive objects.]}

  \item Navas, S.\ et al.\ (Particle Data Group) (2024).
        \textit{Review of Particle Physics}.
        \textit{Physical Review D}, 110, 030001.
        \textit{[Not yet read by the authors. Cited in Section~1.5 for
        the experimental bound on CPT violation from neutral kaon
        systems.]}

  \item Sciama, D.\,W.\ (1953).
        On the origin of inertia.
        \textit{Monthly Notices of the Royal Astronomical Society},
        113, 34--42.
        \textit{[Not yet read by the authors. The Olbers-style
        calculation in Appendix~A is structurally identical to
        Sciama's integral for inertial origin, as noted in that
        appendix.]}

  \item Susskind, L.\ (1995).
        The world as a hologram.
        \textit{Journal of Mathematical Physics}, 36(11), 6377--6396.
        \textit{[Not yet read by the authors. Co-developer of the
        holographic principle with 't~Hooft; cited in Section~1.3
        as a pillar of the argument.]}

  \item Reissner, H.\ (1916).
        \"{U}ber die Eigengravitation des elektrischen Feldes nach der
        Einsteinschen Theorie.
        \textit{Annalen der Physik}, 50(9), 106--120.
        \textit{[Not yet read by the authors. Original derivation of
        the charged spherically symmetric metric via the Einstein
        equations. Section~4 derives the same metric by an
        information-theoretic route, without field equations.]}

  \item Nordstr\"{o}m, G.\ (1918).
        On the energy of the gravitational field in Einstein's theory.
        \textit{Verhandl.\ Koninkl.\ Ned.\ Akad.\ Wetenschap.,
        Afdel.\ Natuurk., Amsterdam}, 26, 1201--1208.
        \textit{[Not yet read by the authors. Independent derivation of
        the same charged metric; the solution bears both authors' names.
        See Reissner (1916) above.]}

\end{enumerate}

%% ============================================================
\appendix
\section{Cosmological Motivation for Assumption A2}
\label{app:A2}
%% ============================================================

Assumption A2 --- that the vacuum near matter operates at holographic
saturation --- is the least self-evident of the four assumptions
underlying the derivation.
We offer here a quantitative motivation, deferring the exact
calculation for the concordance cosmology to Paper~III.

\subsection*{Intellectual context}

The question of where space comes from has a long history.
Leibniz argued against Newton that space has no existence independent
of the relations between objects; to posit an absolute space is to
posit something with no observable consequences and no causal role.
Mach \cite{Mach1883} extended this relational programme to inertia,
proposing that the resistance of matter to acceleration is not a
response to absolute space but to the gravitational influence of all
other matter in the universe.
General relativity absorbed Mach's relational spirit without fully
vindicating it: the metric is relational in form, but the manifold
remains a given substrate.

Sciama (see Other References) gave Mach's proposal its most precise
quantitative form, showing that the sum of gravitational potential
contributions $\sum GM_i/r_i c^2$ from all matter in the observable
universe yields a result of order unity, suggesting that the inertia
of any local object is entirely due to its interaction with
distant matter.
The calculation below is structurally identical to Sciama's integral
(see Other References),
evaluated in information-theoretic language rather than gravitational
potential language.
The agreement is not a coincidence; both are expressions of the same
underlying physics.

\subsection*{The Olbers integral for holographic area}

We model the universe as Einstein--de~Sitter: spatially flat,
matter-dominated, $\Omega = 1$, $\Lambda = 0$, with critical density
$\rho_c = 3H_0^2/8\pi G$ and particle horizon $r_H = 2c/H_0$.
Rather than summing light flux from distant sources as in Olbers's
paradox, we sum holographic area contributions.

From Section~3.3, a mass $M$ at distance $r$ contributes information
in transit
\begin{equation}
  \Delta C = \frac{2\pi Mcr}{\hbar},
  \label{eq:app-transit}
\end{equation}
which by A3 requires holographic area
\begin{equation}
  \Delta A = 4\ell_P^2\,\Delta C = \frac{8\pi GMr}{c^2}.
  \label{eq:app-deltaA}
\end{equation}
Integrating over concentric shells at critical density:
\begin{equation}
  A_{\rm required}
  = \int_0^{r_H}\frac{8\pi Gr}{c^2}\,\rho_c\,4\pi r^2\,dr
  = \frac{32\pi^2 G\rho_c}{c^2}\cdot\frac{r_H^4}{4}
  = 48\pi\frac{c^2}{H_0^2}.
  \label{eq:app-required}
\end{equation}
The holographic capacity of the cosmological horizon is
\begin{equation}
  A_{\rm horizon} = 4\pi r_H^2 = 16\pi\frac{c^2}{H_0^2},
  \label{eq:app-horizon}
\end{equation}
giving the ratio
\begin{equation}
  \frac{A_{\rm required}}{A_{\rm horizon}} = 3.
  \label{eq:app-ratio}
\end{equation}
This is a factor-of-3 result with no free parameters.

By the Copernican principle this holds at every point: each location
receives information flux from its own observable universe, and that
flux saturates the local holographic screen to within a factor of~3
under the Einstein--de~Sitter approximation.
We conjecture that the exact calculation for the concordance cosmology
--- incorporating dark energy and the redshift evolution of matter
density --- closes this gap and yields exact saturation.
The factor of~3 overshoot in the EdS case is qualitatively consistent
with dark energy reducing the effective matter volume relative to
horizon area.

\subsection*{Physical interpretation}

This result provides independent quantitative motivation for A2:
the vacuum operates at holographic saturation not because we
postulate it, but because the collective information flux from all
matter in the observable universe fills the local holographic screen
to capacity at every point.

It also suggests a more complete resolution to the question Leibniz
and Mach were asking.
Space is not a pre-existing stage.
The local vacuum exists, and has the geometric structure it has,
because the universe has been generating information long enough and
densely enough that the past light cone of every point is saturated.
Spacetime is not ordained --- it is accumulated, built incrementally
from the information output of all matter within causal reach.

This reframes the apparent wastefulness of cosmic voids.
Every patch of apparently empty space encodes the complete history
of its past light cone, back to the beginning of time: a holographic
record of everything that has ever been causally accessible from
that location.
Matter is a light dusting of active information-generation centers
embedded in a universe that is overwhelmingly composed of this
accumulated record.

The Machian implication runs deeper still.
The same integral that gives equation~\eqref{eq:app-ratio} is
structurally identical to Sciama's calculation of
inertial origin.
This suggests that mass-energy itself --- and ultimately $E = Mc^2$
--- has a Machian, relational origin: the energy content of matter
is built up from its interactions with the cosmological information
field rather than being intrinsic to matter in isolation.
Screen Theory may thus vindicate Mach's programme in a precise
quantitative sense: both space and inertia are accumulated,
relational quantities, not given substrates.
This connection is developed in Paper~III.

\subsection*{Note on placement}

This appendix presents the argument at the level of a toy calculation
sufficient to establish plausibility.
The full derivation --- with proper cosmological distance measure,
redshift corrections, and exact treatment of the concordance cosmology
--- is reserved for Paper~III, where it forms part of the central
argument.
This appendix may be read independently of the main text and may be
updated or replaced as the Paper~III calculation is completed.

\end{document}
